% ===================================================================
%  तालिका सं. 3 — बचपन और यौवन।
%  Traduction 2026 d'apres texto/io, controlee sur texto/fr et
%  texto/en. Consigne : voir texto/hi/10-talika-01.tex.
%
%  Deux scenes, et l'ido subdivise beaucoup plus que le francais :
%  huit des intertitres n'ont pas de vis-a-vis francais. Le hindi suit
%  l'ido, comme les autres traductions.
% ===================================================================

%%K t03-apar-1 apar
\VUsaut{3.0mm}
\VUtitre{15.0pt}{60}{तालिका सं. 3}
\VUsaut{1.6mm}
\VUtitre{11.4pt}{0}{बचपन और यौवन।}
\VUsaut{3.4mm}

%%K t03-apar-2 apar
(तालिकाएँ 3 और 4 इस प्रकार रची गई हैं कि चार \VUgras{पारिवारिक दृश्यों} में
\VUgras{मौसम}, \VUgras{आयु}, \VUgras{परिवार}, \VUgras{वस्त्र} और
\VUgras{स्वास्थ्य} की मूल शब्दावली सामने आ जाए।)
\VUsaut{4.0mm}

%%K t03-tit-1 sub
\VUcentre{12.0pt}{0}{\textit{पहला दृश्य।}}
\VUsaut{3.0mm}

%%K t03-tit-2 sub
\VUpk{10.2pt}{40}{गाँव में एक बपतिस्मा।}
\VUsaut{2.4mm}

%%K t03-tit-3 sub
\VUpk{10.2pt}{40}{वसंत।}
\VUsaut{3.0mm}

%%K t03-c1-01-1 p
1. --- हम \VUgras{वसंत} में हैं, \VUgras{मार्च}, \VUgras{अप्रैल} या
\VUgras{मई} के महीने में, \VUgras{सप्ताह} के किसी दिन --- \VUgras{सोमवार},
\VUgras{मंगलवार} या \VUgras{बुधवार}। सुबह के \VUgras{दस बजे} हैं।

%%K t03-c1-02-1 p
2. --- \VUgras{मौसम} सुहावना है, \VUgras{हवा} निर्मल। सूरज चमक रहा है। कुछ छोटे-छोटे \VUgras{बादल}
\textsuperscript{(2)} नीले \VUgras{आकाश} \textsuperscript{(1)} में उठ रहे हैं।

%%K t03-c1-03-1 p
3. --- \VUgras{पेड़ों} पर अभी \VUgras{पत्ते} लगभग नहीं हैं। दुबले
\VUgras{चिनार} \textsuperscript{(3)} और ऊँचे \VUgras{प्लेन वृक्ष}
\textsuperscript{(17)} पर अब तक कलियाँ ही हैं।

%%K t03-c1-04-1 p
4. --- \VUgras{अबाबीलें} \textsuperscript{(4)} लौट आई हैं, और आनंद से चीखती
हुई उस \VUgras{शिखर-चूल} \textsuperscript{(7)} के चारों ओर चक्कर काट रही हैं जो
\VUgras{घंटाघर} \textsuperscript{(5)} की है, और वह घंटाघर गाँव के
\VUgras{गिरजे} \textsuperscript{(6)} पर खड़ा है।

%%K t03-tit-4 sub
\VUsaut{3.4mm}
\VUpk{10.2pt}{40}{गिरजा।}
\VUsaut{3.0mm}

%%K t03-c1-05-1 p
5. --- यह गिरजा बहुत सादा है, अपने बड़े \VUgras{द्वार} \textsuperscript{(13)} और
अपनी बड़ी \VUgras{घड़ी} \textsuperscript{(8)} के साथ। \VUgras{छोटी सुई}
\textsuperscript{(9)} X के अंक पर है : वह \VUgras{घंटे} बताती है।
\VUgras{बड़ी वाली} \textsuperscript{(10)} XII (दोपहर) पर है; वह \VUgras{मिनट}
बताती है।

%%K t03-c1-06-1 p
6. --- मीनार में \VUgras{घंटे} \textsuperscript{(11)} प्रसन्नता से बज रहे हैं।
छत के सिरे पर पत्थर का एक छोटा \VUgras{क्रूस} \textsuperscript{(12)} खड़ा है।

%%K t03-c1-07-1 p
7. --- गिरजे में केवल बड़ा \VUgras{द्वार} \textsuperscript{(13)} ही नहीं, बगल में
एक छोटा दरवाज़ा भी है। उसी दरवाज़े से \VUgras{पादरी} \textsuperscript{(15)}, अपने
\VUgras{चोगे} में, \VUgras{वस्त्रागार} \textsuperscript{(14)} से निकलकर
\VUgras{पादरी-निवास} \textsuperscript{(19)} की ओर जाते हैं।

%%K t03-c1-08-1 p
8. --- उनके पास यह रही \VUgras{दूधवाली} \textsuperscript{(24)} अपने
\VUgras{गधे} \textsuperscript{(23)} के साथ। वह शहर से लौट रही है, जहाँ बच्चों के
लिए अपना अच्छा दूध पहुँचा आई है।

%%K t03-tit-5 sub
\VUsaut{4.0mm}
\VUpk{10.2pt}{40}{फलों का बग़ीचा।}
\VUsaut{3.4mm}

%%K t03-c1-09-1 p
9. --- उस्ताद अलीबोरों (गधा) इस सुबह की दौड़ से बहुत प्रसन्न है। वह पड़ोस के
\VUgras{बग़ीचे} के उन \VUgras{फूलों} \textsuperscript{(20)} से महकी हवा आनंद से
भीतर खींचता है जो \VUgras{फलदार पेड़ों} \textsuperscript{(18)} को ढके हैं। वे
\VUgras{आड़ू के पेड़} \textsuperscript{(18)} और \VUgras{खुबानी के पेड़}
\textsuperscript{(19)} बिलकुल गुलाबी और सफ़ेद हैं, और अच्छी फसल की आशा बँधाते
हैं।

%%K t03-c1-10-1 p
10. --- इस बीच नन्ही \VUgras{मधुमक्खियाँ} \textsuperscript{(22)}, जो
\VUgras{छत्ते} \textsuperscript{(21)} से आती हैं, वहाँ जाकर गुनगुनाती हुई अपना
मीठा \VUgras{शहद} बटोरती हैं।

%%K t03-c1-11-1 p
11. --- पर यह क्या हो रहा है? गाँव के बच्चे दौड़ते हुए आ रहे हैं और डरी हुई
\VUgras{बत्तखों} \textsuperscript{(25)} को तितर-बितर कर रहे हैं। यह वह जुलूस है जो
मुंशी के बेटे के \VUgras{बपतिस्मे} के बाद गिरजे से निकल रहा है।

%%K t03-c1-12-1 p
12. --- नन्हा जेम्स, अपने \VUgras{पालने} \textsuperscript{(26)} में, घंटों की
प्रसन्न ध्वनि से जाग उठता है, और उसकी बहन मेडलिन, जो जुलूस देखना चाहती है, अपनी
माँ से कहती है कि आकर देहरी पर ही उसके बाल सँवार दें और उसे कपड़े पहना दें।

%%K t03-c1-13-1 p
13. --- उसकी माँ मान जाती है। वह अपना \VUgras{सिर का कपड़ा} \textsuperscript{(28)},
अपनी \VUgras{चोली} \textsuperscript{(29)} और अपना \VUgras{एप्रन} \textsuperscript{(30)}
पहनकर दरवाज़े के सामने एक छोटी कुर्सी पर आ बैठती है। मेडलिन पर केवल उसका
\VUgras{पेटीकोट} \textsuperscript{(31)} और उसकी \VUgras{कुरती} \textsuperscript{(32)}
है।

%%K t03-c1-14-1 p
14. --- वह कितनी प्रसन्न है! वह अपने प्रिय फूल भूल जाती है — अपने
\VUgras{कार्नेशन} \textsuperscript{(34)}, जिन पर \VUgras{तितलियाँ}
\textsuperscript{(35)} बैठती हैं, अपने \VUgras{ट्यूलिप} \textsuperscript{(36)}, अपनी
\VUgras{घाटी की लिली} \textsuperscript{(37)} जो \VUgras{गमलों}
\textsuperscript{(38)} में हैं, और वे गमले \VUgras{दीवार} \textsuperscript{(33)}
पर रखे हैं, और अपने \VUgras{गुलाब}
\textsuperscript{(39)} जो इतनी सुंदर बाड़ बनाते हैं। वह जुलूस की महिलाओं के बहुमूल्य
वस्त्र निहार रही है।

%%K t03-c1-15-1 p
15. --- गाँव के लड़के वस्त्रों पर कम ध्यान देते हैं। वे \VUgras{मिठाइयाँ}
\textsuperscript{(55)} या \VUgras{सिक्के} \textsuperscript{(52)} पाने के लिए आगे
लपकते और धक्कमधक्का करते हैं। उनमें से एक रो \textsuperscript{(40)} रहा है।

%%K t03-c1-16-1 p
16. --- दूसरे ने \VUgras{कुत्ते} \textsuperscript{(42)} का पंजा कुचल दिया है, जो
कूँ-कूँ करता भाग जाता है और \VUgras{मुर्गी} \textsuperscript{(43)} तथा उसके
\VUgras{चूज़ों} \textsuperscript{(44)} को उड़ा देता है।

%%K t03-tit-6 sub
\VUsaut{5.0mm}
\VUpk{10.2pt}{40}{बपतिस्मे का जुलूस।}
\VUsaut{4.0mm}

%%K t03-c1-17-1 p
17. --- जुलूस के आगे-आगे \VUgras{धाय} \textsuperscript{(45)} चल रही है। उसके सिर
पर लंबे फीतों वाली सुंदर \VUgras{टोपी} \textsuperscript{(46)} है। वह
\VUgras{नवजात शिशु} \textsuperscript{(47)} को लिए है, जो पूरा \VUgras{लेस}
\textsuperscript{(48)} में लिपटा है।

%%K t03-c1-18-1 p
18. --- धाय के पीछे \VUgras{धर्मपिता} \textsuperscript{(49)} और
\VUgras{धर्ममाता} \textsuperscript{(53)} आते हैं। वे मिठाइयाँ और सिक्के बाँट रहे
हैं। धर्मपिता के हाथ में \VUgras{दस्ताने} \textsuperscript{(51)} और सिर पर
\VUgras{ऊँचा हैट} \textsuperscript{(50)} है। धर्ममाता हाथ में एक सुंदर
\VUgras{गुलदस्ता} \textsuperscript{(54)} लिए है।

%%K t03-c1-19-1 p
19. --- उनके पीछे हमें बच्चे के \VUgras{चाचा} \textsuperscript{(56)} और
\VUgras{चाची} \textsuperscript{(58)} दिखाई देते हैं। चाची के सिर पर सजीला
\VUgras{हैट} \textsuperscript{(59)} है, चाचा पर काला \VUgras{लंबा कोट}
\textsuperscript{(57)} और सफ़ेद वास्कट। फिर \VUgras{चचेरा भाई}
\textsuperscript{(60)} और उसकी \VUgras{चचेरी बहन} \textsuperscript{(61)} आते हैं। वह
नवजात शिशु की \VUgras{बहन} \textsuperscript{(62)}, नन्ही बर्था, का हाथ थामे है।

%%K t03-c1-20-1 p
20. --- बर्था का दूसरा \VUgras{भाई} \textsuperscript{(63)} पीछे चल रहा है। वह एक
\VUgras{रिश्तेदार} \textsuperscript{(64)} को हाथ दिए हुए है। परिवार के
\VUgras{मित्र} \textsuperscript{(65)} उनके बाद आते हैं।

%%K t03-tit-7 sub
\VUsaut{4.6mm}
\VUpk{10.2pt}{40}{तमाशबीन।}
\VUsaut{3.4mm}

%%K t03-c1-21-1 p
21. --- जुलूस के पास यह रहा एक \VUgras{भिखारी} \textsuperscript{(66)}, अपनी
\VUgras{बैसाखी} \textsuperscript{(67)} के सहारे। वह भीख माँग रहा है।

%%K t03-c1-22-1 p
22. --- दूसरी ओर एक बहुत \VUgras{शिष्ट} \textsuperscript{(68)} लड़का है। वह अपनी
टोपी उठाकर जान-पहचान वालों को \VUgras{« नमस्ते »} कहता है।

%%K t03-c1-23-1 p
23. --- गाँव वाले बपतिस्मे का जुलूस निकलते देखने अपने घरों से बाहर आ गए हैं।
हमें अपनी \VUgras{सूती टोपी} पहने एक \VUgras{गँवई आदमी} \textsuperscript{(69)}
दिखाई देता है और उसके पास एक \VUgras{गँवई औरत} \textsuperscript{(70)}। फिर एक \VUgras{बुढ़िया}
\textsuperscript{(71)}, जो अपनी \VUgras{छड़ी} \textsuperscript{(72)} के सहारे
है। अंत में \VUgras{चक्कीवाला} \textsuperscript{(73)} अपने चौड़े
\VUgras{नमदे के हैट} \textsuperscript{(74)} में, और
\VUgras{खड़ाऊँ} \textsuperscript{(75)} पहने एक जवान गँवई औरत, जो अपने बच्चे को
चलना सिखा रही है।

%%K t03-c1-24-1 p
24. --- दृश्य बहुत जीवंत है।
\VUsaut{4.0mm}

%%K t03-tit-8 sub
\VUcentre{12.0pt}{0}{\textit{दूसरा दृश्य।}}
\VUsaut{3.0mm}

%%K t03-tit-9 sub
\VUpk{10.2pt}{40}{सार्वजनिक उत्सव।}
\VUsaut{2.4mm}

%%K t03-tit-10 sub
\VUpk{10.2pt}{40}{जल-क्रीड़ा और नौका-दौड़।}
\VUsaut{3.0mm}

%%K t03-c2-01-1 p
1. --- \VUgras{गरमी} आ गई है। \VUgras{जून} (जुलाई या \VUgras{अगस्त}) के महीने
का तपता सूरज नौजवानों को नहाने और नौका चलाने के लिए ललचाता है।

%%K t03-c2-02-1 p
2. --- नदी के दाहिने किनारे पर नाविक जल-क्रीड़ा और नौका-दौड़ में उतरने के लिए
कपड़े उतार रहे हैं।

%%K t03-c2-03-1 p
3. --- उनमें से एक अपने \VUgras{जूते} \textsuperscript{(1)} उतार रहा है; वह उनके
तसमे (बटन) खोल रहा है। उसके दो साथी पहले ही तैयार हैं। अब उन पर केवल अपनी
\VUgras{बनियान} \textsuperscript{(2)} और \VUgras{तैरने का जाँघिया}
\textsuperscript{(3)} है। वे अपने साथियों की प्रतीक्षा करते बातें कर रहे हैं। वे
दूसरे किनारे पर लगे मेले और उत्सव की बात कर रहे हैं।

%%K t03-c2-04-1 p
4. --- उनके पास ज़मीन पर उनके साथियों के \VUgras{कपड़े} पड़े हैं : एक
\VUgras{जोड़ी} \VUgras{बूट} \textsuperscript{(4)}, \VUgras{जूते}
\textsuperscript{(5)}, \VUgras{मोज़े} \textsuperscript{(6)}, \VUgras{नमदे}
\textsuperscript{(7)} और \VUgras{तिनके} \textsuperscript{(8)} के हैट तथा एक
\VUgras{टाई} \textsuperscript{(9)}।

%%K t03-c2-05-1 p
5. --- एक नौजवान ने अपना \VUgras{कोट} \textsuperscript{(10)} सावधानी से तह किया
है। वह उसे एक तौलिये पर रख रहा है, जिसमें उसका पड़ोसी थोड़ी देर में उन दोनों के
\VUgras{सूट} बाँध देगा।

%%K t03-c2-06-1 p
6. --- छठा लड़का पिछड़ गया है। उसके सिर पर अब भी उसकी \VUgras{टोपी}
\textsuperscript{(11)} है। उसने न अपनी \VUgras{कमीज़} \textsuperscript{(12)} उतारी
है, न अपना \VUgras{कॉलर} \textsuperscript{(13)}, न अपने \VUgras{कफ़}
\textsuperscript{(14)}। वह हाथ में अपनी \VUgras{वास्कट} \textsuperscript{(17)} लिए
है और अब भी अपना \VUgras{पतलून} \textsuperscript{(16)} तथा अपनी
\VUgras{गैलिसें} \textsuperscript{(15)} पहने है। वह अगली दौड़ के लिए निश्चय ही
तैयार न हो सकेगा।

%%K t03-c2-07-1 p
7. --- नौजवानों के पास \VUgras{ऊँटकटारों} \textsuperscript{(18)} से घिरी एक पगडंडी
नदी के किनारे उतरती है, जहाँ बूढ़ा जेम्स \VUgras{सरकंडों} \textsuperscript{(19)}
के बीच शाम को छोड़ी जाने वाली \VUgras{आतिशबाज़ी} \textsuperscript{(20)} तैयार कर
रहा है।

%%K t03-tit-11 sub
\VUsaut{4.4mm}
\VUpk{10.2pt}{40}{दौड़।}
\VUsaut{3.4mm}

%%K t03-c2-08-1 p
8. --- नदी पर अपनी छतरियों की छाँह में कुछ महिलाएँ एक \VUgras{डोंगी}
\textsuperscript{(21)} में निकली हैं। वे दौड़ देख रही हैं, जो बहुत रोमांचक है। हर
\VUgras{नाव} \textsuperscript{(22)} में चार \VUgras{खेवैये} \textsuperscript{(24)}
पूरी शक्ति से खींच रहे हैं, जब कि \VUgras{कर्णधार} \textsuperscript{(26)}
\VUgras{पतवार} \textsuperscript{(25)} थामे नाज़ुक नाव को दिशा देता है।
\VUgras{डाँड़} \textsuperscript{(23)} ताल से पानी पर पड़ते हैं, और हलकी नावें
चकरा देने वाली गति से दौड़ती हैं।

%%K t03-c2-09-1 p
9. --- कुछ नौजवान पुल के खंभे के पास \VUgras{चिकने बाँस} \textsuperscript{(28)}
का पुरस्कार जीतने में लगे हैं। उनमें से एक लचकते फिसलन भरे बाँस पर सावधानी से
सरकता हुआ सिरे पर बँधी नन्ही \VUgras{झंडी} \textsuperscript{(29)} तक पहुँचना चाहता
है। उसका अभागा \VUgras{प्रतिद्वंद्वी} \textsuperscript{(30)} पानी में गिर पड़ा है।
वह तैरकर किनारे लौट रहा है।

%%K t03-c2-10-1 p
10. --- बड़े लड़के \VUgras{जल-युद्ध} \textsuperscript{(32)} खेल रहे हैं,
\VUgras{घाट} \textsuperscript{(33)} के पास। एक बड़ी \VUgras{भीड़}
\textsuperscript{(31)} \VUgras{पुल} \textsuperscript{(27)} की \VUgras{मुँडेर} पर
झुकी और \VUgras{किनारे} पर ठसाठस भरी ये सब खेल देख रही है। पर कुछ तमाशबीन मुड़कर
\VUgras{चढ़ने के बाँस} \textsuperscript{(45)} का खेल देखते हैं, या
\VUgras{पुरस्कार-वितरण} की ओर जाते \VUgras{नगराध्यक्ष के जुलूस} को निहारते हैं।

%%K t03-c2-11-1 p
11. --- सबसे पहले यह रहे कुछ \VUgras{छोकरे} \textsuperscript{(35)}, जो
\VUgras{बाजेवालों} \textsuperscript{(36)} के आगे-आगे दौड़ रहे हैं; फिर आते हैं
\VUgras{नगराध्यक्ष} \textsuperscript{(37)}, उनके उप-अध्यक्ष और छोटे नगर की
\VUgras{नगर-परिषद्}। सबसे पीछे \VUgras{दमकल वाले} \textsuperscript{(38)} चलते हैं।

%%K t03-tit-12 sub
\VUsaut{5.0mm}
\VUpk{10.2pt}{40}{मेला।}
\VUsaut{3.6mm}

%%K t03-c2-12-1 p
12. --- \VUgras{मेले के मैदान} \textsuperscript{(39)} पर बड़ी भीड़ है। लोग
तरह-तरह के \VUgras{तंबू} देखने आते हैं : \VUgras{लकड़ी के घोड़े}
\textsuperscript{(40)}, \VUgras{सरकस} \textsuperscript{(41)}, जहाँ खेल आरंभ भी हो
चुका है; \VUgras{सार्वजनिक नाच} \textsuperscript{(42)}, जहाँ नगर के नौजवान लड़के
और लड़कियाँ \VUgras{नाचने} का सुख ले रहे हैं।

%%K t03-c2-13-1 p
13. --- सिरे पर हमें \VUgras{अखाड़ा} \textsuperscript{(44)} दिखाई देता है, जहाँ
वीर स्पेनी \VUgras{साँड़-योद्धा} क्रुद्ध \VUgras{साँड़} \textsuperscript{(45)} से
लड़ता है; फिर \VUgras{झूला-रेल} \textsuperscript{(49)} और
\VUgras{निशानेबाज़ी का तंबू} \textsuperscript{(46)}, जहाँ लोग \VUgras{बंदूक}
चलाने और \VUgras{निशाने} पर मारने का अभ्यास करते हैं।

%%K t03-c2-14-1 p
14. --- पास ही यह रहा \VUgras{मेले का रंगमंच} \textsuperscript{(47)}, जहाँ
\VUgras{प्रहसन} और \VUgras{नाटक} खेले जाते हैं। उसके बिलकुल पास
\VUgras{दानव} \textsuperscript{(48)} और \VUgras{बौने} का तंबू है। पहले की
\VUgras{ऊँचाई} दो मीटर पंद्रह है; दूसरा मुश्किल से एक बच्चे जितना ऊँचा है।

%%K t03-c2-15-1 p
15. --- अंत में यह रही बड़ी \VUgras{चिड़ियाघर-गाड़ी} \textsuperscript{(50)}, अपने
\VUgras{जंगली जानवरों} के साथ — अपने \VUgras{बाघ} \textsuperscript{(51)} के, जिसके
\VUgras{पंजे} बड़े बलवान हैं, अपने \VUgras{सिंह} \textsuperscript{(52)} के, जिसकी
\VUgras{अयाल} घनी है, अपनी दो \VUgras{जिराफ़ों} \textsuperscript{(53)} और अपने
\VUgras{हाथी} \textsuperscript{(54)} के, जो अपनी \VUgras{सूँड़} इतनी चतुराई से काम
में लाता है।

%%K t03-c2-16-1 p
16. --- \VUgras{जानवर सधाने वाला} \textsuperscript{(55)}, जो इन जानवरों को
सिखाता है, तंबू के दरवाज़े पर खड़ा है।
\VUsaut{6.0mm}
\VUfilet{22mm}
